
\section*{\underline{EXPOSÉ VIII}\\[3pt] \underline{\underline{LE THÉORÈME DE RIEMANN--ROCH}}}
\begin{center}
\underline{par Pierre BERTHELOT}
\end{center}

\subsection*{1. Morphismes d'intersection complète}

\textbf{Définition 1.1.} \uline{Soit} @@M0@@ \uline{un morphisme de schémas}. \uline{On dit que} @@M1@@ \uline{est un morphisme d'intersection complète si pour tout point} @@M2@@ \uline{de} @@M3@@, \uline{il existe un voisinage ouvert} @@M4@@ \uline{de} @@M5@@, \uline{un} @@M6@@-\uline{schéma lisse} @@M7@@, \uline{tels que la restriction de} @@M8@@ \uline{à} @@M9@@ \uline{se factorise en} :
@@M10@@
\uline{où} @@M11@@ \uline{est une immersion régulière} (VII 1.4).

Quitte à remplacer @@M12@@ par un ouvert, on peut supposer que @@M13@@ est une immersion fermée régulière.

On peut dans cette définition remplacer @@M14@@ par un fibré vectoriel @@M15@@. En effet, la notion étant locale, on peut supposer que @@M16@@ est affine, et que son image dans @@M17@@ est contenue dans un ouvert affine @@M18@@ de @@M19@@. Il existe alors un entier @@M20@@ et une @@M21@@-immersion fermée @@M22@@, d'où une @@M23@@-immersion @@M24@@, qui est régulière d'après VII 1.10. Comme la composée de deux immersions régulières est régulière, on trouve donc une @@M25@@-immersion régulière @@M26@@ factorisant @@M27@@.

\textbf{Proposition 1.2.} \uline{Soit} @@M28@@ \uline{un morphisme d'intersection complète}. \uline{Alors toute} @@M29@@-\uline{immersion} @@M30@@ \uline{de} @@M31@@ \uline{dans un} @@M32@@-\uline{schéma lisse est régulière}.

Nous montrerons d'abord le corollaire :

\textbf{Corollaire 1.3.} \uline{Soit le diagramme commutatif} :
@@M33@@
